\section{Introduction}
\label{sec:intro}

The emergence of large-scale pre-trained models, such as CLIP \cite{radford2021learning} and GPT-3 \cite{brown2020language}, has revolutionized machine learning. To adapt these models to downstream applications, parameter-efficient fine-tuning and task adaptation have become standard practices. However, fine-tuning independently on diverse downstream tasks often results in isolated, single-task expert models. In multi-task scenarios, deploying multiple specialized backbones incurs severe storage and inference costs. 

To address this, model merging has emerged as a promising paradigm for multi-task learning without joint retraining. By directly interpolating or combining the weights of independent, fine-tuned models, model merging aims to construct a single unified backbone that retains multi-task capabilities \cite{wortsman2022model, choshen2022fusing}. The foundational method, Task Arithmetic \cite{ilharco2022editing}, creates task vectors by subtracting the pre-trained weights from the fine-tuned weights and scaling them with a single global scalar coefficient per task. 

While Task Arithmetic is simple and effective, weight-space interference often degrades performance when merging diverse or conflicting tasks. To resolve this, recent state-of-the-art (SOTA) model merging frameworks have shifted toward increasingly complex, fine-grained, and parameter-rich merging schemes. Prominent among these are layer-wise test-time adaptive merging methods like AdaMerging \cite{yang2024adamerging} and SyMerge \cite{jung2025symerge}. These methods optimize distinct merging coefficients for each transformer layer and each task, under the assumption that different layers of the model capture distinct semantic information and thus contribute differently to task performance. Under the banner of Test-Time Adaptation (TTA), these coefficients are tuned on unlabeled target data, often using entropy minimization or self-labeling objectives.

\begin{figure}[t]
\centering
\includegraphics[width=\columnwidth]{results/fig1_treatments.png}
\caption{Model merging average accuracies across different diagnostic treatments on CLIP ViT-B/32, aggregated over 3 independent trials. Our multi-seed, dual-optimizer analysis exposes a profound interaction between optimizer choice, task complexity, and parameter regularization, showing that layer-specificity is an illusion under zero-order search but emerges as a delicate, overfit configuration under unconstrained first-order gradient descent.}
\label{fig:treatments_intro}
\end{figure}

In this paper, we take a highly critical, methodology-focused sanity check of this foundational assumption. As \textbf{The Methodologist}, our core philosophy is that bad experimental practices, weak baselines, and unexamined assumptions lead to false progress. We initially identify three major potential methodological gaps in current model merging literature:
\begin{enumerate}
    \item \textbf{The Layer-Specificity Assumption:} Recent papers claim that optimized layer-wise coefficients represent fine-grained, layer-specific representational contributions across tasks. However, this assumption has never been rigorously validated against diagnostic control treatments.
    \item \textbf{Under-tuned and Weak Baselines:} SOTA publications frequently compare their complex, multi-parameter test-time adaptation methods against a weak, manually selected uniform Task Arithmetic baseline. They fail to benchmark against properly tuned, single-scalar task-wise baselines.
    \item \textbf{Landscape Flatness Ignored:} The optimization landscape of merging coefficients is rarely characterized. If the landscape is extremely flat, learned parameter variations may simply be optimization noise acting as a coarse regularizer, rather than meaningful, precise, layer-specific allocations.
\end{enumerate}

To expose and analyze these issues, we construct a rigorous \textbf{Sanity-Checking and Interpretability Suite} to evaluate optimized merging coefficients on CLIP ViT-B/32 across four diverse vision classification tasks (MNIST, FashionMNIST, CIFAR-10, SVHN) across \textbf{3 independent random seeds}. We implement three main diagnostic treatments on learned coefficients: (1) \textbf{Intra-Task Layer Shuffling} (Shuffle Treatment), (2) \textbf{Task-Wise Spatial Averaging} (Mean Treatment), and (3) \textbf{Norm-Bounded Perturbation} (Noise Treatment). Furthermore, we conduct representational similarity analysis using linear Centered Kernel Alignment (CKA) \cite{kornblith2019similarity} to evaluate activation alignment with original expert models. Crucially, we isolate optimizer confounding by deploying two distinct optimization regimes: a zero-order \textbf{Adaptive 1+1 Evolution Strategy (1+1 ES)} and a first-order backpropagation-based \textbf{Adam Gradient Descent (Adam GD)}.

Our empirical results yield a profound and unexpected dialectical finding, which we term the \textbf{Overfitting-Optimizer Paradox}:
\begin{itemize}
    \item \textbf{Layer-Specificity is an Illusion under Zero-Order Search:} Under zero-order 1+1 ES search, layer-specificity behaves like an illusion. Replacing complex, layer-wise optimized coefficients (52 parameters) with their flat, task-wise Spatial Mean (4 parameters, a 92.3\% reduction) actually improves average test accuracy from $85.07 \pm 0.47\%$ to $85.21 \pm 0.11\%$. Here, Spatial Averaging acts as a powerful regularizer that eliminates high-frequency zero-order mutation noise.
    \item \textbf{The Overfitting Paradox and Delicate Layer-Specificity under Adam GD:} Under first-order Adam GD, unconstrained optimization finds a highly precise, delicate configuration of layer coefficients that is extremely sensitive to shuffling (causing a $5.43\%$ drop in average performance and a $15.69\%$ collapse on CIFAR-10) or spatial averaging (causing a $10.35\%$ collapse on CIFAR-10), creating an illusion of physical layer-specificity. However, this delicate structure fails to outperform the unoptimized Task Arithmetic baseline ($84.44 \pm 0.37\%$) on the unseen test set while introducing 4x greater seed variance ($1.57\%$ vs. $0.37\%$), confirming that its layer-specificity is a delicate transductive overfitting artifact on the small calibration set.
    \item \textbf{Extreme Landscape Flatness:} Both optimizers tolerate up to \textbf{50\% relative Gaussian noise} on coefficients with negligible average decay (retaining $83.89 \pm 0.32\%$ accuracy for 1+1 ES), confirming that the merging optimization landscape resides in an exceptionally flat basin.
    \item \textbf{CKA vs. Downstream Accuracy Discrepancy:} While spatial averaging improves activation similarity on average ($+0.0069$ CKA for 1+1 ES and $+0.0036$ CKA for Adam GD), high-level kernel alignment is a poor predictor of downstream classification performance because activation correlation decouples from weight-space decision boundary integrity, with CIFAR-10 accuracy collapsing by over 10\% under averaging despite maintaining high activation correlation ($>0.95$ CKA).
\end{itemize}

This study serves as a critical, nuanced course-correction for the model merging community. We demonstrate that unconstrained test-time adaptation of layer-wise merging coefficients is highly prone to transductive overfitting. While zero-order 1+1 ES overfitting manifests as high-frequency optimization noise easily smoothed out by spatial averaging (which acts as a regularizer), first-order Adam GD finds a highly delicate, overfit configuration that collapses under shuffling or averaging, despite achieving no generalizable performance gains over the unoptimized baseline. We urge the community to adopt rigorous optimizer controls, task-complexity-aware diagnostic baselines, and explicit coefficient regularization to prevent this overparameterized transductive drift.
